%% Appendix H — OTA Procedures
\chapter{OTA Update Procedures}
\label{app:ota}

\section{OTA Architecture}

MMS V2 uses the ESP-IDF \texttt{esp\_https\_ota()} function for firmware updates. The dual-partition scheme allows automatic rollback if the new firmware is unstable.

\begin{figure}[H]
\centering
\begin{tikzpicture}
\matrix[matrix of nodes,
    nodes={draw=mmsdark!70, minimum width=58mm, minimum height=7mm, anchor=center, font=\small},
    row sep=-\pgflinewidth] (flash) {
Bootloader \hfill 4KB\\
Partition table \hfill 3KB\\
NVS \hfill 24KB\\
OTA data \hfill 8KB\\
app0 firmware slot \hfill 1MB\\
app1 firmware slot \hfill 1MB\\
LittleFS logs/config \hfill 2MB\\
};
\node[right=8mm of flash-3-1, font=\scriptsize, align=left] {config survives OTA};
\node[right=8mm of flash-4-1, font=\scriptsize, align=left] {tracks active partition};
\node[right=8mm of flash-5-1, font=\scriptsize, align=left] {currently booted slot};
\node[right=8mm of flash-6-1, font=\scriptsize, align=left] {OTA writes alternate slot};
\node[right=8mm of flash-7-1, font=\scriptsize, align=left] {offline sessions and revocations};
\end{tikzpicture}
\caption{ESP32 4MB flash layout for OTA}
\label{fig:ota_flash_layout}
\end{figure}

\section{OTA Update Procedure}

\subsection{Step 1: Build the Firmware}

In Arduino IDE:
\begin{enumerate}[leftmargin=1.5em]
    \item Open \texttt{v2/access\_node/access\_node.ino}
    \item Update version: in \texttt{config.h}, change \texttt{\#define FW\_VERSION "2.0.0"} to the new version
    \item Verify compile: Sketch $\rightarrow$ Verify/Compile (Ctrl+R)
    \item Export binary: Sketch $\rightarrow$ Export Compiled Binary
    \item The \texttt{.bin} file appears in the sketch directory
\end{enumerate}

\subsection{Step 2: Upload to Firebase Storage}

\begin{enumerate}[leftmargin=1.5em]
    \item Firebase Console $\rightarrow$ Storage $\rightarrow$ \texttt{firmware/} folder
    \item Click Upload File, select the \texttt{.bin} file
    \item Name it by version: \texttt{2.1.0.bin}
    \item After upload: click the file $\rightarrow$ three-dot menu $\rightarrow$ Copy download URL
    \item The URL is a long HTTPS URL; copy it exactly
\end{enumerate}

\subsection{Step 3: Trigger OTA}

Option A — Via web app:
\begin{itemize}[leftmargin=1.5em]
    \item Super Admin Panel $\rightarrow$ OTA Update $\rightarrow$ Paste URL and version $\rightarrow$ Submit
\end{itemize}

Option B — Via Firebase Console (RTDB):
\begin{enumerate}[leftmargin=1.5em]
    \item Open RTDB $\rightarrow$ navigate to \texttt{/mms/ota}
    \item Set the value to:
\begin{lstlisting}[style=jsonlstyle]
{
  "firmware_url":   "https://firebasestorage.googleapis.com/.../2.1.0.bin",
  "target_version": "2.1.0",
  "released_at":    1718600000
}
\end{lstlisting}
\end{enumerate}

Option C — Via Firebase CLI:
\begin{lstlisting}[style=bashstyle]
firebase database:update /mms/ota --data '{
  "firmware_url":   "https://...",
  "target_version": "2.1.0",
  "released_at":    1718600000
}'
\end{lstlisting}

\subsection{Step 4: Monitor Progress}

On the dashboard: each node card shows ``updating...'' during OTA and reverts to ``idle'' with the new version number on success.

On Serial Monitor (if USB connected):
\begin{verbatim}
[OTA] Received OTA command: v2.0.0 -> v2.1.0
[OTA] URL: https://firebasestorage.googleapis.com/...
[RELAY] OFF (safety before OTA)
[OTA] Starting download...
[OTA] Progress: 10%
[OTA] Progress: 25%
...
[OTA] Progress: 100%
[OTA] Verification OK
[OTA] Rebooting to new firmware...
--- REBOOT ---
MMS V2 Access Node
FW: 2.1.0
...
[NVS] fw_version updated: 2.1.0
[MQTT] Published status: {fw_version: "2.1.0", ...}
\end{verbatim}

\section{OTA Rollback}

The firmware includes automatic rollback:
\begin{itemize}[leftmargin=1.5em]
    \item If the new firmware crashes 3 times on boot, the bootloader reverts to the previous partition
    \item The node reports the old firmware version on reconnect
    \item Manual rollback is not possible via OTA; re-trigger OTA with the previous known-good \texttt{.bin} file
\end{itemize}

\section{Manual Flash Recovery}

If OTA fails catastrophically (rollback also fails):
\begin{enumerate}[leftmargin=1.5em]
    \item Connect USB cable from laptop to ESP32
    \item Open Arduino IDE with the last known-good sketch
    \item Upload via USB (\textit{not} OTA)
    \item After successful upload: node boots from the new firmware on app0 partition
    \item Re-run integration tests before returning node to service
\end{enumerate}

\section{OTA Safety Checklist}

Before triggering OTA on production nodes:
\begin{itemize}[leftmargin=1.5em]
    \item[ ] New firmware tested on bench node (direct USB flash, then integration tests 1--8)
    \item[ ] OTA tested on one bench node before pushing to all production nodes
    \item[ ] Last known-good \texttt{.bin} file archived locally
    \item[ ] Relay turns OFF before OTA starts (safety: machines don't stay powered during update)
    \item[ ] RTDB OTA \texttt{released\_at} timestamp is current Unix time (not a stale value)
    \item[ ] All nodes showing ``idle'' before OTA (no active sessions)
\end{itemize}
